\subsection{The unaveraged point comparison}
\paragraph{The actual weighted coefficient comparison.}
\begin{proposition}[The unaveraged weighted theta comparison]
\label{ex:app-theta-comparison}
Put $m=|h_*|$ and $F=\mathbf Q(\sqrt{Dh_*})$, where $D$ is a
positive squarefree parameter in the fixed odd filter. Assume
$(D/\ell_0)=1$ at the chosen actual stage, so that $\ell_0$ is
inert in $F$. Keep the local Schwartz functions, the actual function
$P_M(z)=\lambda_M(L\pi_{\ell_0}(z))$, and the integer $c_0$
of (O.4). Then a $2$-power root of unity $\zeta_M(D)$ satisfies
\[
 H_M(mD)=c_0\,2^{\omega(N)}\zeta_M(D)
                   \lambda_M(\overline{LP_D})\pmod {2^M}.\tag{O.7}
\]
The root of unity belongs to the field generated by the fixed
characters at the primes of $N$. No class number is divided out.
\end{proposition}


\begin{proof}
The filter makes $Dh_*$ an odd fundamental discriminant, with
$F$ split at every prime dividing $2N$. Its unit group is
$\{\pm1\}$. The condition $Q_B(x)=mD$ is equivalent to
$x^2=Dh_*$. Together with $(1+x)/2\in R_{z,2}$ and the other
integrality conditions, this says exactly that
\[
 \mathcal O_F\longrightarrow R_z,
 \qquad \sqrt{Dh_*}\longmapsto x
\]
is an optimal embedding. Optimality follows because the source
order is already maximal. All nonzero local weights on such an
embedding are roots of unity. The valuation-one condition at each
prime of $h_*$ holds because $(D,h_*)=1$; at $\ell_0$ the residue
of $x$ is nonzero. The fixed unit filters hold by admissibility.

First classify the local embeddings. Away from $N\ell_0$ there
is one local orbit: an $\mathcal O_{F,p}$-stable lattice is locally
a fractional ideal, hence principal, and this also covers ramified
primes dividing $Dh_*$. At $p^e\Vert N$, the algebra is split.
The two integral idempotents decompose both lattices of the
Eichler lattice pair. Its quotient is cyclic of order $p^e$,
so its elementary exponents lie wholly in one summand. The two
possibilities give exactly two orientations. This applies at two
using the integral spectral projectors specified in the construction.
At $\ell_0$, the maximal order of the division algebra has residue
field $\mathbf F_{\ell_0^2}$. The two residue-field embeddings
of the unramified quadratic order give the two orientations.
Consequently there are $2^{\omega(N)+1}$ local orientation tuples.

Fix one tuple. Choose a rational embedding
$\iota:F\hookrightarrow B$ and write $x_0=\iota(\sqrt{Dh_*})$.
Such an embedding exists because $F$ is a field at the two ramified
places of $B$, namely $\ell_0$ and infinity
\cite[Proposition 14.6.7, $(v)\Rightarrow(iii)$]{voight-quaternion}.
Choose $z_0$ realizing
the selected integral local embeddings. The adelic representatives
with this tuple are exactly
\[
       \widehat F^{\times}z_0\widehat R^{\times}.
\]
Their quotient by the rational centralizer $F^{\times}$ is
\[
 F^{\times}\backslash\widehat F^{\times}/
      (\widehat F^{\times}\cap z_0\widehat R^{\times}z_0^{-1})
 =F^{\times}\backslash\widehat F^{\times}/\widehat{\mathcal O}_F^{\times}
 =\operatorname{Pic}(\mathcal O_F).
\]
Thus this is an actual class-group torsor, not an average or a
quotient of a class-group orbit.

For an ideal-class representative $a\in\widehat F^{\times}$,
the element $a$ centralizes $x_0$, and therefore
\[
 f((az_0)^{-1}x_0(az_0))=f(z_0^{-1}x_0z_0).
\]
On the other hand,
\[
 \chi_{h_*}(\operatorname{Nrd}(az_0))
 =\chi_{h_*}(\operatorname{N}_{F/\mathbf Q}a)
                    \chi_{h_*}(\operatorname{Nrd}z_0).
\]
The first factor is exactly the genus character in (O.1).
To see that it is well-defined on the class group, for $D>1$
the biquadratic field $\mathbf Q(\sqrt D,\sqrt{h_*})$ has
discriminant $D\,h_*\,(Dh_*)=(Dh_*)^2$. Its relative
discriminant over $F$ is therefore one. The character obtained
by restricting $\chi_{h_*}$ to $G_F$ is unramified and, under
class field theory, is $a\mapsto
\chi_{h_*}(\operatorname{N}_{F/\mathbf Q}a)$. For $D=1$ it
is the trivial character. Ideal representatives can be chosen
prime to $h_*$, giving exactly the elementary norm formula in
(O.1). Reversing the Artin convention has no effect because this
character is quadratic.

